%!TEX root=../fade.tex

This section gives a rough breakdown of the \fade{} code: how to install, compile and run it.

You will need an internet connection, and some knowledge of how to use a terminal environment (or a well configured \href{https://visualstudio.microsoft.com/}{Visual Studio} installation). The steps here will work on any POSIX compatible system (Linux or MacOS) -- on windows machines you must either \href{https://learn.microsoft.com/en-us/windows/wsl/install}{install WSL}, or have a working \texttt{msvc} compiler installed on your machine.

\subsection{Installing the Code}

	The code is hosted on github at \url{https://github.com/DrFraserGovil/FADE-Emulator/}. It can either be downloaded manually, or cloned in the terminal:

	\begin{terminal}
>> git clone https://github.com/DrFraserGovil/FADE-Emulator
	\end{terminal}

	\def\jsl{\textmd{\textsc{jsl}}}

\subsection{Compiling the Code}
	At the present, we do not distribute precompiled binaries, so it is necessary to compile the code on your own machine. You must first \href{https://cmake.org/download/}{install cmake} and ensure that you have a valid compiler (any C++20 compatible compiler is fine).

	Then:
	\begin{terminal}
>> cd FADE-Emulator
>> make
	\end{terminal}

	You must have an internet connection when you run this command, as it will download\footnote{Using CMake FetchContent} some other software: the \href{https://libeigen.gitlab.io/}{Eigen linear algebra library}\footnote{If you are familiar with python, this is \texttt{numpy.linalg} on steroids}, and the \href{https://github.com/DrFraserGovil/JSL}{Jack Standard Library} (\jsl{}), a separate project I have authored that handles much of the `application' code.

	It will spit out a bunch of text  - as of 03/08/26, these include a number of warnings (in purple) - but it should complete the compilation.

	\subsubsection{Validating Compilation}

		Running the command
		\begin{terminal}
			./fade test
		\end{terminal}

		Places the code into testing mode. If you have a valid python environment (with \texttt{matplotlib} installed), it will generate a graphic demonstrating the $W_i$ fractions for a randomly instantiated model.

		If you don't have Python - you will get some error messages from your OS, but a \texttt{test} directory will have been creating, holding the associated data.
\subsection{\fade{} and \jsl{}}

	The internal workings of the \fade{} make heavy use of an existing library I have written, the \jsl{}.

	\subsubsection{Logging}
		\fade{} uses the \href{https://jack-standard-library.readthedocs.io/en/stable/docfiles/Log.html#log}{custom \jsl{} logger} to write the output that you see when writing it. This allows some of the output to be suppressed when you want the code to be `quiet', or output much more data when you want to know what has gone wrong.

		\begin{enumerate}
			\item \texttt{./fade \ldots} gives the normal output level
			\item \texttt{./fade --quiet \ldots} suppresses all output except error messages
			\item \texttt{./fade --verbose \ldots} produces lots of output which can be used for diagnostics, but it unsuitable for general use
		\end{enumerate}

	\subsubsection{Configuration}

		The FADE Emulator has a large number of parameters that can be modified. You may view them by typing:

		\begin{terminal}
			./fade help
		\end{terminal}

		This activates the help menu, and details which values which can be changed by the user, and a description of what they do. Most settings have a similar name to the names used in this document (though not always).

		These settings are grouped into named groups, but this has no meaning for the user, other than helping to group parameters that do similar things together.

		Changing the value is done in one of two ways:

		\paragraph{Command Line Arguments} Running the \texttt{./fade} executable and adding extra arguments after it:

		\begin{terminal}
./fade train -expert 3,3 -model my_file.fde -file training_data.dat
		\end{terminal}
		This will attempt to launch the training module with:
		\begin{enumerate}
			\item Only 3 experts (the range $3 \to 3$, see below for multi-model training syntax)
			\item Beginning the training at the previous model defined by \texttt{my\_file.fde}
			\item Using the training data stored in \texttt{training\_data.dat}
		\end{enumerate}
		All other values will use their default value.

		Note that some parameters have multiple aliases which can be used: \texttt{-input} and \texttt{file} have the same meaning. This is detailed in the \texttt{help} menu.

		\paragraph{Configuration files} As you tweak more and more values, the number of commands given to run the code will get longer and longer. Instead, you can write your data to a config file:

		\begin{terminal}
// test.cfg
expert 3,3
model my_file.fde
file training_data.dat
(...)

>> ./fade --config test.cfg
		\end{terminal}

		This loads the data in the config file, and reads it as if it were passed as arguments to the command line.


		\paragraph{Mixed-Mode} It is possible to use both a config file and command line arguments:

		\begin{terminal}
		>> ./fade --config test.cfg -model other_model.fde
		\end{terminal}

		In this case, the config file is loaded \textit{first}, and then the command line values are applied, overwriting any values previously set: in the example above, the value passed to \texttt{model} is `\texttt{other\_model.fde}': the one given at the command line and not `\texttt{my\_file.fde}' (the one in the config file).


		\begin{aside}[{To dash or not two dash?}]
			When specifying command line variables, any number of dashes are allowed: \texttt{--flag} and \texttt{-flag} are always equivalent. The general advice is to use one dash for short commands (\texttt{-f}) and two for long ones (\texttt{--model}), but this is not a requirement.

			When specifying variables in a config-file, do not use leading dashes.
		\end{aside}
\subsection{Training Models}

	In order to create a new model, the user must provide \textbf{training data}. The expected data format is as follows:

	\begin{terminal}
x_0 x_1 x_2 (...) x_n z y
// repeat for each y on next lines
	\end{terminal}
	I.e., a space-delimited enumeration of the position vector $\x_t$, the prior-weighting $\zeta_t^a$ and the observed value $y_t^a$.

	If $\x_t$ is one dimensional ($\x_t = x_t$), the data file would look like:

	\begin{terminal}
//training.data
x_1 z_1 y_1
x_2 z_2 y_2
x_3 z_3 y_3
(...)
\end{terminal}

	If you do not know what $\zeta_t^a$ is (\ref{E:DefineWeights}), then set these values to 1 (but do not omit them entirely).

	\begin{aside}[Deprecation Imminent]
		This method of inputing data will be deprecated soon, in favour of one which natively supports $\x$-clustering. At the moment, clustering is achieved automatically by grouping values of $\x$ which are closer than the \texttt{ClusteringRadius} - a user-configured variable.
	\end{aside}

	Once the training data has been generated, the model may be trained:

	\begin{terminal}
		./fade train --file training.data --save my_model.fde --expert 3,6 --dep 2,4
	\end{terminal}

	This will generate a model called \texttt{my\_model.fde} which has been trained on the training data. This will simultaneously train the set of submodels for $3 \leq N_e \leq 6$ and $2 \leq N_d \leq 4$ (for a total of 12 individual models).
\subsection{Using Models}

	Once a model has been trained, it can be used to make predictions of the distribution function $p(y| \x)$. In order to achieve this, the user must write a file with the following syntax:

	\begin{terminal}
//qfile.dat
x_0 x_1 x_2 (...) x_n y_1 y_2 y_3 y_4 y_5 (...) y_m
//repeat for each value of X
\end{terminal}

	As with the training data, the \texttt{x\_0 ...x\_n} specify the values of $\x$, a point in emulation space to be queried. The values of $y$ form the grid that the model will be evaluated on: so specifying \texttt{y\_0 y\_1 y\_2 ..} means ``at position $\x$, tell me the probability density at $y_0$, then $y_1$, then $y_2$.'' There is no requirement that $\{y_i\}$ be sorted or uniform.

	Multiple $\x$-queries can be entered, with each new $\x-y$ prediction entered on a new line.

	The prediction is then made by calling:

	\begin{terminal}
./fade predict --query qfile.dat --model my_model.fde
	\end{terminal}

	The output of a prediction is a second data file (with name configured by \texttt{query-out}). If the model which was queried was trained on multiple $N_e, N_d$, then a separate query file (denoted by \texttt{[name]\_Nd\_Ne.dat}) is written for each pair.

	For each $\vec{x}-y$ specification in the query file, the output file takes the form:

	\begin{terminal}
New query: x_0 x_1 (...) x_n
y_1 y_2 y_3 (...) y_m
p_1 p_2 p_3 (...) p_m
//repeat for each x-y in query
\end{terminal}
	Where:
	\begin{equation}
		p_i = p^\text{cns}(y_i | \x)
	\end{equation}
